\documentclass[a4paper, 12pt, left=30mm, right=15mm, top=20mm, bottom=20mm]{report}
\usepackage[utf8]{inputenc}
\usepackage[russian]{babel}
\usepackage{mhchem} % for 13C
\usepackage{verbatim} % for multiline commentary 
\usepackage{graphicx} % for logo
\usepackage{titling} % for shift margins
\usepackage{blindtext} % for text-generation
\usepackage{setspace} % to set interval
\usepackage[numbers,sort&compress]{natbib} % for range cite [1-4]

\begin{document}

\begin{titlepage}

	\begin{centering}
		\includegraphics[width=0.25\textwidth]{logo.png}\par
	\end{centering}
	\centerline{Московский государственный университет имени М. В. Ломоносова}
	\centerline{Факультет вычислительной математики и кибернетики}
	\centerline{Кафедра математической кибернетики}
	\centerline{\hfill\hrulefill\hrulefill\hfill}
	\vfill
	\vfill
	\large
	\centerline{Стешин Семен Сергеевич}
	\vfill
	\Large
	\begin{centering}
		\textbf{Khnum: быстрая open-source программа \\ для расчета метаболических потоков \\ с использованием \ce{^{13}C}-углерода}
		
	\end{centering}
	\normalsize
	\vfill
	\centerline{Выпускная квалификационная работа}
	\vfill
	\vfill
	\begin{flushright}
	Научный руководитель:\\
	к.ф.м.н., доцент \\
	Шуплецов М. С.
	\end{flushright}
	\vfill
	\vfill
	\centerline{Москва --- 2020}

\end{titlepage}



\begin{abstract}
В биологии и медицине встречается задача определения скорости метаболических потоков внутри клетки. 
Один из методов решения этой задачи --- \ce{^{13}C}-Metabolic Flux Analysis --- анализ метаболических потоков с использованием  \ce{^{13}C}-углерода. 
В этом методе, исследователи проводят эксперимент и обрабатывают его результаты на компьютере. Проблема в том, что современные программы для анализа метаболических потоков либо имеют закрытый код и платны для коммерческого использования, либо написаны неэффективно, из-за чего вычисления могут занимать недели для одного эксперимента. В этой работе проведен краткий обзор метода, написана эффективная программа для решения задачи и проведено сравнение с существующими аналогами.
\end{abstract}

\tableofcontents

\chapter{Введение}
\section{Мотивация}
Рак --- вторая по частоте причина смерти в мире\cite{Cancer_statistics}. Сто лет назад Отто Варбург заметил\cite{Warburg_effect} особенность раковых клеток: они склонны производить энергию с помощью активного гликолиза, вместо более эффективного окислительного фосфорилирования. Знание этого позволило находить опухоли с помощью позитронно-эмиссионной томографии, а Варбурга наградили Нобелевской премией.

Диабетом болеет 8.8\% людей в мире\cite{Diabetes_statistics}. Почти 4 миллиона в год умирает из-за этой болезни. Лечения пока нет, но есть симптоматическая терапия инъекциями инсулина. Раньше его получали из поджелудочных желез свиней и коров, но препарат было сложно очистить, поэтому иногда случались аллергические реакции. Все изменилось в 1978 году, когда компания Genentech смогла создать генетически-модифицированную кишечную палочку, которая в ходе жизнедеятельности производила чистый человеческий инсулин\cite{Genentech_paper}. Сейчас таким образом производят почти весь препарат.

В случае с эффектом Варбурга, открытие заключалось в изменении скорости химической реакции, протекающей внутри клетки. В случае с инсулином, решается задача метаболической инженерии --- увеличить скорость синтеза инсулина, не убив кишечную палочку. В обоих случаях надо уметь измерять скорости внутриклеточных химических реакций -- их называют потоками. Один из современных методов измерения потоков -- \emph{\ce{^{13}C}-Metabolic Flux Analysis} (далее \emph{MFA}), что переводится как анализ метаболических потоков. Этому методу посвящена наша работа.

MFA имеет множество применений в исследованиях рака\cite{Application_cancer_2009, Application_cancer_2012, Application_cancer_2013, Application_cancer_2015, Application_cancer_2017, Application_cancer_2018, Application_cancer_2018_2}
, в метаболической инженерии\cite{Application_engeneering_2009, Application_engeneering_2015, Application_engeneering_2017} и в других областях\cite{Application_other_2011, Application_other_2013, Application_other_2014}.

\section{Эксперимент}
\section{Компьютерное моделирование}

\chapter{Основные понятия}
\section{Глоссарий}
\section{Прямая задача}
\section{Обратная задача}
\chapter{Постановка задачи}
\chapter{Основная часть}
\chapter{Полученные результаты}


\begin{thebibliography}{XXXX}
	\bibitem{Cancer_statistics}
	Всемирная Ассоциация Здравоохранения. Cancer [Электронный ресурс] URL: https://www.who.int/news-room/fact-sheets/detail/cancer (дата обращения: 12.03.2020)
	
	\bibitem{Warburg_effect}
	Warburg O., Wind F., Negelein E. The metabolism of tumors in the body //The Journal of general physiology.--- 1927. --- Т. 8. --- №. 6. --- С. 519.
	
	\bibitem{Diabetes_statistics}
	Zimmet P. et al. Diabetes mellitus statistics on prevalence and mortality: facts and fallacies //Nature Reviews Endocrinology. --- 2016. --- Т. 12. --- №. 10. --- С. 616.
	
	\bibitem{Genentech_paper}
	Cohen S. N. et al. Construction of biologically functional bacterial plasmids in vitro //Proceedings of the National Academy of Sciences. --- 1973. --- Т. 70. --- №. 11. --- С. 3240--3244.
	
	\bibitem{Application_cancer_2009}
	Metallo C. M., Walther J. L., Stephanopoulos G. Evaluation of 13C isotopic tracers for metabolic flux analysis in mammalian cells //Journal of biotechnology. --- 2009. --- Т. 144. --- №. 3. --- С. 167--174.
	
	\bibitem{Application_cancer_2012}
	Walther J. L. et al. Optimization of 13C isotopic tracers for metabolic flux analysis in mammalian cells //Metabolic engineering. --- 2012. --- Т. 14. --- №. 2. --- С. 162--171.
	
	\bibitem{Application_cancer_2013}
	Hiller K., Metallo C. M. Profiling metabolic networks to study cancer metabolism //Current opinion in biotechnology. --- 2013. --- Т. 24. --- №. 1. --- С. 60--68.
	
	\bibitem{Application_cancer_2015}
	Boroughs L. K., DeBerardinis R. J. Metabolic pathways promoting cancer cell survival and growth //Nature cell biology. --- 2015. --- Т. 17. --- №. 4. --- С. 351--359.
	
	\bibitem{Application_cancer_2017}
	Dong W., Keibler M. A., Stephanopoulos G. Review of metabolic pathways activated in cancer cells as determined through isotopic labeling and network analysis //Metabolic engineering. --- 2017. --- Т. 43. --- С. 113--124.
	
	\bibitem{Application_cancer_2018}
	Antoniewicz M. R. A guide to 13 C metabolic flux analysis for the cancer biologist //Experimental \& molecular medicine. --- 2018. --- Т. 50. --- №. 4. --- С. 1--13.
	
	\bibitem{Application_cancer_2018_2}
	Badur M. G., Metallo C. M. Reverse engineering the cancer metabolic network using flux analysis to understand drivers of human disease //Metabolic engineering. --- 2018. --- Т. 45. --- С. 95--108.
	
	\bibitem{Application_engeneering_2009}
	Nakahigashi K. et al. Systematic phenome analysis of Escherichia coli multiple‐knockout mutants reveals hidden reactions in central carbon metabolism //Molecular systems biology. --- 2009. --- Т. 5. --- №. 1.
	
	\bibitem{Application_engeneering_2015}
	Crown S. B., Long C. P., Antoniewicz M. R. Integrated 13C-metabolic flux analysis of 14 parallel labeling experiments in Escherichia coli //Metabolic engineering. --- 2015. --- Т. 28. --- С. 151--158.
	
	\bibitem{Application_engeneering_2017}
	Long C. P. et al. Enzyme I facilitates reverse flux from pyruvate to phosphoenolpyruvate in Escherichia coli //Nature communications. --- 2017. --- Т. 8. --- №. 1. --- С. 1--8.
	
	\bibitem{Application_other_2011}
	Wahrheit J., Nicolae A., Heinzle E. Eukaryotic metabolism: measuring compartment fluxes //Biotechnology journal. --- 2011. --- Т. 6. --- №. 9. --- С. 1071--1085.
	
	\bibitem{Application_other_2013}
	Metallo C. M., Vander Heiden M. G. Understanding metabolic regulation and its influence on cell physiology //Molecular cell. --- 2013. --- Т. 49. --- №. 3. --- С. 388--398.
	
	\bibitem{Application_other_2014}
	Dieuaide-Noubhani M., Alonso A. P. (ed.). Plant metabolic flux analysis: methods and protocols. --- Humana Press, 2014.
	
\end{thebibliography}

\end{document}          
